% SPDX-License-Identifier: Apache-2.0
%
% TxHDL, step by step: a tutorial for somebody meeting the language for
% the first time. Built by //docs:tutorial. Every listing is included
% from the tree, and every printout and figure is what the build made
% by running the example, so the tutorial cannot say what the code
% does not do. Two columns, as the paper is; the listings are set at
% the paper's size, and the waveforms are drawn at column width.
\documentclass[journal,9pt]{IEEEtran}

\newcommand{\listingsize}{\fontsize{6}{7}\selectfont}
% SPDX-License-Identifier: Apache-2.0
%
% The house preamble, shared by every document in this directory. Loaded
% after \documentclass, before \title. Keeping it in one file is what
% stops two articles drifting apart in font, listing style or figure
% style.
% Computer Modern. IEEEtran selects Times (ptm) by default, so the face has
% to be chosen explicitly.
%
% Latin Modern is the Computer Modern variety used here, and the choice is
% forced rather than aesthetic. Under T1 encoding, plain `cmr` has no Type 1
% outlines unless cm-super is installed, and it is not in the pinned
% distribution, so pdflatex falls back to EC bitmaps and every glyph in the
% PDF becomes a Type 3 font. That was measured: the first build produced a
% document with 10 Type 3 fonts and no Type 1. Latin Modern is Computer
% Modern redrawn with a full T1 repertoire and real .pfb outlines, so the
% same design comes out scalable.
\usepackage[T1]{fontenc}
\usepackage[utf8]{inputenc}
\usepackage{lmodern}
% microtype is not in the pinned TeX distribution. emergencystretch alone
% keeps the two-column measure from producing overfull lines.
\setlength{\emergencystretch}{3em}

\usepackage{amsmath}
\usepackage{amssymb}
\usepackage{amsthm}
\usepackage{booktabs}
\usepackage{array}
% enumitem is not in the pinned TeX distribution; the list spacing below
% does what \setlist would have done.
\usepackage{float}
\usepackage{xcolor}
\usepackage{listings}
\usepackage{tikz}
\usetikzlibrary{arrows.meta,positioning,fit,backgrounds,calc}
\usepackage[hidelinks,bookmarks=true]{hyperref}

% Numbered, definition-style box for a rule the reader is meant to apply.
\theoremstyle{definition}
\newtheorem{rulebox}{Rule}
\newtheorem{example}{Example}

% Unnumbered asides.
\theoremstyle{remark}
\newtheorem*{tip}{Tip}
\newtheorem*{pitfall}{Pitfall}

% Tighter lists, in place of enumitem's \setlist.
\setlength{\topsep}{2pt}
\setlength{\itemsep}{1pt}
\setlength{\parsep}{0pt}

\newcommand{\code}[1]{\texttt{#1}}
\newcommand{\term}[1]{\emph{#1}}

% Syntax highlighting. Four roles, distinguishable in colour and still
% distinguishable in greyscale, because an IEEE article gets printed: the
% keyword is the only bold role, the comment is the only italic one, and the
% two remaining roles differ in lightness as well as in hue.
\definecolor{kwblue}{RGB}{20,60,150}
\definecolor{cmtgreen}{RGB}{90,120,90}
\definecolor{strred}{RGB}{150,50,40}
\definecolor{numgray}{gray}{0.45}
\definecolor{framegray}{gray}{0.70}
\definecolor{bgshade}{gray}{0.97}
% A darker fill for figure cells. The listing background has to stay very
% light to keep code readable; a figure cell has to be clearly shaded or the
% distinction it draws is invisible in print.
\definecolor{cellshade}{gray}{0.90}

% The listing size is the document's choice, because it depends on the
% column: a two-column article sets `\listingsize` to 6pt and keeps its
% listings within 70 columns, a one-column document keeps the body size
% and includes source files at 80. Lines are never broken; a line that does not fit
% overflows the frame, which is visible, rather than wrapping, which is
% not.
\providecommand{\listingsize}{\scriptsize}

% `'` is deliberately not a string delimiter: the language uses it for loop
% labels, as Rust does, so treating it as a quote would swallow the rest of
% the line.
\lstdefinestyle{txhdl}{
  basicstyle=\ttfamily\listingsize,
  keywordstyle=\bfseries\color{kwblue},
  commentstyle=\itshape\color{cmtgreen},
  stringstyle=\color{strred},
  numberstyle=\tiny\color{numgray},
  morekeywords={mod,use,pub,const,type,struct,enum,trait,impl,fn,async,let,mut,
    match,if,else,loop,while,for,in,break,continue,return,as,await,
    module,bus,channel,transaction,protocol,pipeline,flow,seq,config,
    port,stage,pipe,instance,connect,bind,tag,domain,blackbox,foreign,
    property,role,signal,handler,true,false,
    sequence,interface,modport,implements,package,import,var,wire,func,proc,
    case,default,next,fork,branch,join,state,fsm},
  morecomment=[l]{//},
  morestring=[b]",
  showstringspaces=false,
  columns=fullflexible,
  keepspaces=true,
  breaklines=false,
  % Framed, so a listing is bounded rather than floating in the column.
  frame=single,
  rulecolor=\color{framegray},
  framesep=3pt,
  backgroundcolor=\color{bgshade},
  xleftmargin=3pt,
  xrightmargin=3pt,
  aboveskip=6pt,
  belowskip=6pt,
  % Every listing is numbered and captioned, so it can be referred to by
  % number. `caption` without `float` numbers the listing and leaves it
  % where it was written, which is what a listing in running prose needs.
  captionpos=b,
  abovecaptionskip=2pt,
  belowcaptionskip=6pt,
}

% Figures drawn as text. Framed the same way, and with the highlighting
% switched off, because a box-and-arrow diagram is not code and half its
% words turning blue reads as an accident. Setting a style adds keys rather
% than clearing the ones already set, so the three roles are neutralised by
% hand rather than by leaving them out.
\lstdefinestyle{ascii}{
  basicstyle=\ttfamily\listingsize,
  keywordstyle=\ttfamily,
  commentstyle=\ttfamily,
  stringstyle=\ttfamily,
  columns=fullflexible,
  keepspaces=true,
  frame=single,
  rulecolor=\color{framegray},
  framesep=3pt,
  backgroundcolor=\color{bgshade},
  xleftmargin=3pt,
  xrightmargin=3pt,
  aboveskip=6pt,
  belowskip=6pt,
}
\lstset{style=txhdl}

% House TikZ styles. `shorten` lives on the arrow styles rather than on each
% arrow, so every arrow in the document gets the gap, including ones added
% later. TikZ clips a path at a node's border, so without it an arrowhead
% butts against the box with no gap at all. Keep `node distance` well above
% twice the shorten, or the arrow shrinks to nothing.
\tikzset{
  box/.style={draw,rounded corners=2pt,align=center,inner sep=4pt,
              font=\footnotesize},
  boxf/.style={box,fill=bgshade},
  lbl/.style={font=\scriptsize\itshape,align=center},
  fl/.style={-{Stealth[length=4pt]},shorten >=2.5pt,shorten <=2.5pt},
  fld/.style={fl,dashed},
}



\InputIfFileExists{buildstamp.tex}{}{}
\providecommand{\buildstamp}{Draft build (outside the Bazel flow).}

\title{TxHDL, Step by Step}

\author{Filip Filmar and Dragiša Janković%
\thanks{Both authors are independent. Filip Filmar is the
corresponding author, at f@hdlfactory.com; Dragiša Janković is
reached at linkedin.com/in/dragisa-jankovic.}%
\thanks{This is the tutorial: for somebody who has written some Rust,
has not written hardware, and wants to write a unit, run it, see its
waveform, and get a netlist out, in that order. Every listing is
included from the project repository \code{HDL/txhdl} at build time
and every printout is what the build produced by running it. The
language itself is stated in the embedding article~\cite{embedding},
the runtime in its own document~\cite{runtime}, every example in the
examples document~\cite{examples}, and the cover~\cite{cover} lists
all of them.}%
\thanks{The authors used a large language model, Claude, as an
assistant in exploring the concepts and in writing and constructing
the documents and the programs; every commit of the repository says
so and carries the prompts that led to it, verbatim.}%
\thanks{\buildstamp}}

\markboth{TxHDL, Step by Step}{Filmar and Janković: TxHDL, Step by Step}

\begin{document}
\maketitle

\begin{abstract}
TxHDL is a hardware description language that is a Rust library. A
design is a struct whose state is registers, its behaviour an
\code{async fn} that waits for a clock edge, reads, and drives; the
same file simulates under Rust, writes a waveform, and lowers to
Verilog and VHDL that the build checks against the simulation. This
tutorial takes you from nothing to that, in six steps, each on one
small example the build runs: a counter, the model of time, the three
ways of choosing, channels between units, the waveform, and the
netlist with its checks. It ends with how to add an example of your
own, which is how everything here was added.
\end{abstract}

\begin{IEEEkeywords}
Hardware description languages, Rust, tutorial, simulation, Verilog,
VHDL.
\end{IEEEkeywords}

\section{Before You Start}
\label{sec:t-start}

You need one thing installed: \code{bazelisk}, which fetches the
right Bazel. Everything else, the Rust toolchain, the simulators, the
\LaTeX{} that made this page, comes through the build. Clone the
repository and build it:

\begin{lstlisting}[style=ascii]
bazel build //...
bazel test  //...
\end{lstlisting}

The first build fetches toolchains and takes a while; the second is
seconds. If both finish, you have everything this tutorial uses.

Three sentences on what you are about to read. A \emph{unit} is a
struct of registers with an \code{async fn run} that is its process.
A process \emph{waits} for an edge of a clock, \emph{reads} its
registers and inputs, and \emph{drives} its registers and outputs; the
drives land at the end of the step, so the whole cycle reads as
plain Rust with one rule about time. And a unit written in a subset
the lowering reads is also a netlist, from the same file.

\section{Step One: A Counter}
\label{sec:t-counter}

Listing~\ref{lst:t-counter} is the whole of a unit. It counts to
eight while enabled and raises \code{tick} on the last count. Read it
top to bottom.

\lstinputlisting[rangebeginprefix=//\ begin\{,rangebeginsuffix=\},rangeendprefix=//\ end\{,rangeendsuffix=\},includerangemarker=false,linerange=unit-unit,caption={The counter: a register, a wait, two reads, a predicated drive and an output. From \code{ex\_cheat.rs}.},label={lst:t-counter}]{lib/examples/ex_cheat.rs}

The struct holds the state: one register, \code{n}, three bits wide.
\code{Reg<T>} is a register of any value type; \code{U<3>} is an
unsigned of three bits. \code{\#[derive(Trace)]} lets a waveform find
the field by its name, and \code{Default} gives every register its
zero.

The \code{impl} names the unit's ports in its type: one input, a
\code{Bit} called \code{enable}, and one output, a \code{Bit} called
\code{tick}. \code{In} and \code{Out} are the two ends of a wire; the
testbench holds the other ends. \code{\#[lower]} on the \code{impl}
is what makes the file a netlist as well; leave it off and the unit
still simulates.

The body is one \code{loop}, and the first line of the loop is the
wait: \code{DefaultClock::rising().await}. Everything after it is one
cycle. \code{get} reads the input as the edge left it, and
\code{self.n} in an expression is the register as the edge left it.
\code{with!} is the drives of the unit: under \code{enable},
\code{n} becomes \code{n + 1} at the end of the step, the arrow
\code{<=} pointing into the struct and the entry written as a
struct literal writes a field.
\code{tick.set} drives the output, this cycle, from the value the
register had at the top. Nothing here branches, and nothing is
assigned twice; that is the shape every unit has.

Now run it. \code{bazel run //lib/examples:ex\_cheat} runs the
\code{main} below the unit, which you will read in
Section~\ref{sec:t-testbench}, and prints Listing~\ref{lst:t-counter-out}:
the counter's Verilog, which is Step Six. Fig.~\ref{fig:t-counter}
is its waveform, drawn by the build from the run: \code{n} counts
while \code{enable} is high, \code{tick} marks the sevens, and both
stop when the enable drops.

\begin{figure}[htbp]
\centering
\input{paper_cheat_timing.tex}
\caption{The counter, from its trace: \code{enable} drops at cycle
10, \code{n} counts while it is high, \code{tick} marks the sevens.}
\label{fig:t-counter}
\end{figure}

\lstinputlisting[style=ascii,caption={What \code{ex\_cheat} printed: the counter's Verilog, from the same file.},label={lst:t-counter-out}]{out_cheat.txt}

\section{Step Two: Time}
\label{sec:t-time}

Everything you need to know about time fits on this page, and the
figure above shows it. Time is in \emph{ticks}. The default clock has
a period of two ticks and rises at the even ones, so a cycle is two
ticks and every rising edge is an even number on the axis.

A \emph{register} takes its value at an edge. When the loop drives
\code{self.n <= n + 1}, nothing changes yet; the drive is remembered,
and at the end of the step every drive lands at once. So a register
read after the wait is the value as the edge left it, and a register
driven after the wait shows the new value in the next cycle. That is
why \code{tick} is computed from \code{n} before the drive: the
output in a cycle is a function of the state at the start of it,
which is what hardware does, and what the netlist will do.

A \emph{wire} is what \code{set} drives: an output, or a \code{let}
inside the loop, which the lowering keeps as a named wire. A wire has
no memory; it is a function of the state and the inputs of this
cycle.

The one rule: \textbf{every \code{.await} is a cycle boundary.} A
process that has crossed an edge at this step cannot cross it again,
so a second wait written after the first waits for the next edge. If
you want two events at one edge, you write them in
\code{parallel!}, and the examples document has one that does. State
read before a wait, at the top of the loop, is state read in the
previous cycle; read after the wait, always.

\section{Step Three: Choosing}
\label{sec:t-choosing}

Hardware has no branch. Both arms of a choice exist at once and a
condition selects, which is why the three ways of choosing have their
own names, so that the file reads as the hardware does. \code{with!}
you have seen: the drives of the unit, each under a condition if it
has one, a group of them under one with an optional \code{else}.
\code{when!(c => self \{ .. \})} is one such group with the
condition written first, for a line that is about the condition.
\code{case!} is the group with many arms, and it is what a state
machine is made of. Rust's own \code{if} is there too,
over a \code{bool}, with the drives under it written as \code{set};
its arms are the same priority chain, read as written, and they nest. Listing~\ref{lst:t-case} is a
sequencer that waits for \code{go}, loads a count, runs it down and
reports done.

\lstinputlisting[caption={A state machine with \code{case!}: each arm is a Rust pattern, guards included, and the first that matches wins. \code{ex\_case.rs}.},label={lst:t-case}]{lib/examples/ex_case.rs}

The arms are Rust patterns. An arm with a guard, \code{State::Run if
n == 0}, sits above the plain arm for the same state, and the first
arm that matches wins, exactly as \code{match} would. The state is an
enum of your own: \code{\#[derive(Value)]} gives it a width, two
bits for four variants, and a look in a waveform.
Fig.~\ref{fig:t-case} shows it running: \code{go} for one cycle,
then the states one per edge, and \code{count} running down. The
printout, Listing~\ref{lst:t-case-out}, is the same run as text, and
then the state machine's Verilog: every arm is there, in order, and
nothing else.

\begin{figure}[htbp]
\centering
\input{paper_case_timing.tex}
\caption{The sequencer: \code{go} for a cycle, then \code{Load},
three cycles of \code{Run}, \code{Done}, and back to \code{Idle}.}
\label{fig:t-case}
\end{figure}

\lstinputlisting[style=ascii,caption={What \code{ex\_case} printed: one state per cycle, then the lowering.},label={lst:t-case-out}]{out_case.txt}

The third way is \code{select!}, which chooses a \emph{value} rather
than a drive: \code{select!(f3.raw() => \{ 0 => a + b, 1 => a << b,
.. \})} is a multiplexer, and the Vreteno core's ALU is one of
these. \code{mux(c, a, b)} is the two-way form. You choose between
values with \code{select!} and \code{mux}, between drives with
\code{with!}, \code{when!}, \code{case!} and \code{if}, and that is
the whole vocabulary of choice.

\section{Step Four: Talking Between Units}
\label{sec:t-channels}

Two units talk over a \emph{channel}. A channel carries a
\emph{transaction}, a value with the trait \code{Transaction}, which
any \code{U<N>} has and a struct of values gets by deriving it. It
has two ends: \code{Tx}, on which a unit \code{send}s, and \code{Rx},
on which a unit \code{recv}s. Between them the runtime keeps an
elastic buffer of two, so a send lands in the next cycle and a
receiver never sees a word before the edge after it was offered.
Listing~\ref{lst:t-stage} is a unit between two channels: it waits
\code{until} a word is offered and there is room to pass it on, then
receives it, sends it on incremented, and counts.

\lstinputlisting[rangebeginprefix=//\ begin\{,rangebeginsuffix=\},rangeendprefix=//\ end\{,rangeendsuffix=\},includerangemarker=false,linerange=unit-unit,caption={A stage between two channels: \code{until} is the wait, and its condition guards everything after it. From \code{ex\_stage.rs}.},label={lst:t-stage}]{lib/examples/ex_stage.rs}

\code{until(clock, condition)} is a wait like \code{rising}, one that
passes only at an edge at which the condition holds. \code{peek} is
the head of the channel as the edge left it, \code{ready} whether
there is room on the other channel, \code{recv} the take, \code{take}
the take of whatever is offered as \code{(offered, value)}, for a
unit that serves the channel every cycle, and \code{send} the offer. In the netlist a channel is a \code{valid}, a \code{ready}
and the data, and the lowering makes the wait's condition the guard
of every drive, the \code{ready} of the receive and the \code{valid}
of the send. Fig.~\ref{fig:t-stage} shows the handshake on both
sides of the stage.

\begin{figure}[htbp]
\centering
\input{paper_stage_timing.tex}
\caption{The stage: a word offered on \code{inp} every other cycle,
taken, and passed on incremented on \code{out} a cycle later.}
\label{fig:t-stage}
\end{figure}

A unit that runs every cycle, as a processor core does, cannot wait
\code{until}. It talks by predicate instead: it receives with no wait
before the receive, which takes whatever is offered, and sends under
\code{if}, whose condition becomes the channel's \code{valid}.
Listing~\ref{lst:t-tap} is that form: a tap that takes every word and
passes on the odd ones.

\lstinputlisting[rangebeginprefix=//\ begin\{,rangebeginsuffix=\},rangeendprefix=//\ end\{,rangeendsuffix=\},includerangemarker=false,linerange=unit-unit,caption={A tap: a receive with no wait, and a send under \code{when!}. From \code{ex\_tap.rs}.},label={lst:t-tap}]{lib/examples/ex_tap.rs}

One thing to know before you wire two units with plain wires rather
than a channel: a wire between units carries no promise about which
of the two sees the other's value in the same step. A channel does.
When you need a wire between units, drive it from a register and run
the unit that drives it first; the Vreteno document has the one
case where that was needed, and the reasons.

\section{Step Five: Running and Watching}
\label{sec:t-testbench}

The \code{main} of an example is its testbench. Listing~\ref{lst:t-main}
is the counter's: it makes the two wires with \code{signal}, keeping
the ends the unit does not take; starts a waveform if the build asked
for one; runs the unit with \code{Running::new}; and then drives
\code{enable} and steps the design one cycle at a time with
\code{sim.cycle()}, which is where a testbench looks at wires and
registers between steps.

\lstinputlisting[rangebeginprefix=//\ begin\{,rangebeginsuffix=\},rangeendprefix=//\ end\{,rangeendsuffix=\},includerangemarker=false,linerange=main-main,caption={The counter's testbench: wires, a waveform, the unit running, twelve cycles. From \code{ex\_cheat.rs}.},label={lst:t-main}]{lib/examples/ex_cheat.rs}

The waveform is the part you will use most. \code{Wave::from\_env()}
returns a writer when \code{TXHDL\_FST} or \code{TXHDL\_VCD} names a
file, so the same binary runs silent or writes a trace, and the build
sets the variable when it draws a figure. You name what to watch:
the clock, each wire by a name of your choosing, and a unit by name,
which brings every field it derives \code{Trace} for, under that
name. Then \code{start()}, and every step from there on is written
until \code{stop()}. Open the file in any waveform viewer, or let the
build draw it: every figure in this tutorial is a
\code{waveform()} entry in \code{docs/BUILD.bazel} that names the
example and the signals, and the build runs the example, converts the
trace and draws the picture, so a figure can never disagree with the
code.

A value of your own, an enum or a struct with \code{\#[derive(Value)]},
shows in a trace by its variant name or field by field, as
\code{state} did in Fig.~\ref{fig:t-case}.

One more thing you will want to watch: a signal inside the unit that
is neither a register nor a port. A \code{let} in the loop lives
only in the step, so the trace cannot see it. Keep it as a field of
type \code{Wire<T>} instead, drive it with \code{set} where the
\code{let} was, and read it back with \code{get} in the same step:
it shows in the trace under the unit's name, and in the netlist it is
exactly the wire the \code{let} would have been, so the hardware
does not change. The divider in the examples document~\cite{examples}
keeps its two decisions that way, and the Vreteno core keeps its
stall, its interrupt and its redirect so.

A unit runs in a plain \code{\#[test]} too, with the same
\code{signal}, \code{Running::new} and \code{cycle} as in a
\code{main}. The runtime's clock and executor are per thread, so the
tests of a crate run in parallel as Rust runs them, each with a time
of its own. Compare the unit's registers after a cycle, not its
output wires, which a process sets during the step from the state
before it. A \code{\#[test]} is the place for assertions about
values; the netlist check needs a trace, and comes from the
\code{waveform()} entry of Step Seven.

\section{Step Six: From Rust to Hardware}
\label{sec:t-netlist}

You have seen the counter's Verilog in Listing~\ref{lst:t-counter-out}.
It came from \code{Counter::verilog("counter")}, and
\code{Counter::lowered("counter")} is the netlist itself, which the
runtime can also write as VHDL. Both come from the same \code{run}
you read in Step One: \code{\#[lower]} reads the loop, and the
lowering rewrites what it finds. A register is a register. A wire
named by \code{let} is a wire of that name. \code{with!}, \code{when!},
\code{case!} and \code{if} are predicated assignments in the clocked
block.
\code{select!} and \code{mux} are conditional expressions. The wait
is the clock edge of the block, and an \code{until} is a guard on
everything after it. A memory, \code{Mem<T, N>}, is an array,
written through \code{at} and read by index.

That is the subset: a loop of one wait, then reads, \code{let}s that
are functions of what was read, and drives. Inside it you have
Rust's operators on the value types, \code{+}, \code{-}, \code{*},
\code{\%}, unary \code{-}, \code{\&}, \code{|}, \code{\^{}}, \code{!},
\code{<<}, \code{>>}, the compares and \code{as} to an unsigned
integer, then \code{slice}, \code{concat}, \code{sext} and
\code{zext} for taking words apart and putting them together, and
\code{mul::<M>} for a product wider than its operands. An \code{if}
that yields a value and a \code{match} lower as \code{mux} and
\code{select!} do, and a range, as a pattern or in
\code{contains}, as the two compares it stands for. What you do not
have inside a lowered loop is what hardware does not have: a function
call the lowering does not know, or a loop with a variable count.
The macro says so at compile time, with a message that names the
construct.

The lowering is checked, not trusted. For every lowered example the
build runs the Rust simulation, keeps its trace, writes the VHDL and
the Verilog, and generates a testbench from the trace that replays
the inputs into the netlist and asserts every register and every
output at every tick; then it simulates the VHDL under nvc and the
Verilog under Verilator, both built from source by the build. The
tests are \code{//docs:cheat\_sim\_counter\_tb\_test} and
\code{//docs:cheat\_vsim\_test} for the counter, and there is a pair
for every unit in this tutorial. When you write a unit that lowers,
the same two tests appear for it the moment you add it to the build,
which is Step Seven.

\section{Step Seven: Adding Your Own}
\label{sec:t-adding}

Every example you have read was added the same way, and it is the
way to add yours.

\begin{enumerate}
\item Write \code{lib/examples/ex\_\emph{name}.rs}: the unit, and a
  \code{main} that runs it, prints what is worth reading, and, if
  the unit lowers, writes its netlist with
  \code{txhdl::netlist::write\_vhdl\_from\_env} and prints its
  Verilog. Copy \code{ex\_cheat.rs} and change it.
\item Add a \code{rust\_binary} for it in
  \code{lib/examples/BUILD.bazel}, next to the others, and add the
  file to the \code{sources} filegroup there. Now \code{bazel run
  //lib/examples:ex\_\emph{name}} runs it.
\item Add a \code{waveform()} entry in \code{docs/BUILD.bazel}
  naming the example and the signals to draw; give it
  \code{lowered = (entity, unit)} if the unit lowers. That makes the
  figure, the printout, and the two netlist tests.
\item Write a section for it in \code{docs/examples\_sections},
  include it from \code{docs/examples.tex}, and add the run and the
  figure to the examples document's \code{data}. Build the document
  and look at the page.
\end{enumerate}

The rule behind the four steps is in the repository's
\code{AGENTS.md}: a concept lives in three places, the runtime, an
example, and a document, or it is not finished. The example is the
test.

\section{Where Next}
\label{sec:t-next}

The examples document~\cite{examples} has thirty of these, one per
concept, in the order of the language, with the runtime's public
interface at the end for looking things up. The embedding
article~\cite{embedding} states the language: the rule that every
construct Rust provides is deleted, and what is left. The runtime
document~\cite{runtime} is the library itself, verbatim, with the
model of time stated in prose, for when you want to know what a
construct does inside. And the Vreteno document is a RISC-V core
written in the subset you have just read, checked every cycle against
a reference model, lowered, and placed and routed at 100~MHz, which
is what the six steps add up to.

\begin{thebibliography}{10}

\bibitem{embedding}
F.~Filmar and D.~Janković, ``Deleting a language: Embedding TxHDL in
Rust,'' project repository \code{HDL/txhdl}, \code{//docs:embedding},
2026.

\bibitem{runtime}
F.~Filmar and D.~Janković, ``The TxHDL runtime,'' project repository
\code{HDL/txhdl}, \code{//docs:runtime}, 2026.

\bibitem{examples}
F.~Filmar and D.~Janković, ``TxHDL by example,'' project repository
\code{HDL/txhdl}, \code{//docs:examples}, 2026.

\bibitem{cover}
F.~Filmar and D.~Janković, ``TxHDL documents,'' project repository
\code{HDL/txhdl}, \code{//docs:cover}, 2026.

\end{thebibliography}

\end{document}
